% ===================================================================
%  preambule.tex — macros communes aux deux transcriptions.
%
%  Ce fichier ne contient AUCUNE mesure. Toutes les mesures sont dans
%  kalibro-io.tex et kalibro-fr.tex, un fichier par fac-simile : les
%  deux livrets ne sont pas du meme format, ne sont pas composes au
%  meme corps, et n'ont pas ete numerises de la meme maniere. Charger
%  la calibration AVANT ce fichier ; celui-ci s'en sert et ne la
%  redefinit jamais.
%
%  La methode est celle du projet « Kompleta Gramatiko » :
%  UNE LIGNE DU FAC-SIMILE = UNE LIGNE DU PDF, meme pagination, memes
%  folios, memes coupures de mot. Aucune cesure n'est decidee par TeX ;
%  toutes sont relevees sur le scan et ecrites dans la source.
% ===================================================================

% -------------------------------------------------------------------
% 1. PAQUETAGES
% -------------------------------------------------------------------
\usepackage[T1]{fontenc}
\usepackage[utf8]{inputenc}
\usepackage{XCharter}
\usepackage{geometry}
\usepackage{graphicx}
\usepackage{tikz}
\usepackage{calc}
% L'EXPANSION DE FONTE, pour combler ce que le corps ne comble pas.
% Meme au plus grand corps qui tienne, une page garde des lignes plus
% courtes que les autres : le corps est borne par la ligne LA PLUS
% LONGUE de la page, et les autres s'etirent. Le compositeur de 1926
% reglait cela a l'interlettrage -- il ouvrait ses lignes creuses.
% L'equivalent moderne est l'expansion : chaque glyphe s'etire ou se
% resserre de 2,5 % au plus, ce qui absorbe le manque sans toucher aux
% blancs inter-mots, donc sans les rendre beants. Les coupures de ligne
% ne bougent pas -- elles sont toutes forcees par le releve.
\usepackage[tracking=true,letterspace=0,
            expansion=true,stretch=25,shrink=25]{microtype}

\geometry{
  paperwidth=\VUpapierLargeur, paperheight=\VUpapierHauteur,
  textwidth=\VUtexteLargeur,   textheight=\VUtexteHauteur,
  top=\VUmargeSup, includehead=false, includefoot=false,
  hcentering, % les deux livrets ont des marges laterales egales
}

% -------------------------------------------------------------------
% 2. COUPURES DE LIGNE EXPLICITES
% -------------------------------------------------------------------
% \nl : fin de ligne SANS trait d'union.
% \cc : fin de ligne AVEC trait d'union (le trait est compose).
% On n'utilise jamais \\ : il annule la justification de la ligne.
\newcommand{\nl}{\penalty-10000\relax}
\newcommand{\cc}{-\penalty-10000\relax}

% \parplein TERMINE un alinea dont la derniere ligne court jusqu'a la
% marge — cas d'un alinea qui se poursuit sur la page suivante.
% \parfillskip est GLOBAL : pose une fois, il justifierait la derniere
% ligne de tous les alineas suivants. On l'enferme dans un groupe.
\newcommand{\parplein}{{\parfillskip=0pt\relax\par}}
% \ccplein : la page finit sur un MOT COUPE, et l'alinea continue au
% feuillet suivant. C'est \cc et \parplein a la fois, et ni l'un ni
% l'autre ne fait l'affaire seul :
%   -- \cc pose « -\penalty-10000 ». En fin d'alinea, cette penalite
%      ouvre une derniere ligne VIDE, qui prend un rang de plus sur la
%      page : la page composee gagne une ligne que le fac-simile n'a
%      pas, et tout ce qui suit descend d'un cran.
%   -- \parplein seul ne compose pas le trait d'union.
% On ecrit donc le trait, puis on termine l'alinea a la marge, sans
% penalite. Trois pages du livret ido finissent ainsi (feuillets 66,
% 103, 104) ; le controle 5 les avait signalees comme « \VUcontinue
% sans \parplein », ce qui etait le symptome et non la cause.
\newcommand{\ccplein}{-{\parfillskip=0pt\relax\par}}

% \VUcontinue OUVRE un alinea qui continue la page precedente : pas de
% renfoncement, la ligne reprend au fer a gauche.
\newcommand{\VUcontinue}{\noindent}

% LES ESPACES-MOTS SONT INSECABLES.
% Sans cela le modele fuit en silence : \nl force bien une coupure a la
% fin de la ligne, mais n'empeche pas TeX d'en placer une AVANT si la
% matiere ne tient pas. Une ligne 1 % trop large se recouperait donc
% toute seule, sans meme produire d'« Overfull \hbox ». Les espaces
% rendus insecables, le seul point de coupure legal est celui que le
% releve a place : une ligne trop large devient un « Overfull \hbox »,
% que le controle 1 inventorie. La faute passe du silence au bruit.
% Les espaces gardent toute leur elasticite : la justification est
% intacte.
\makeatletter\newcommand{\VUespaceMot}{%
  \ifdim\spaceskip=0pt\relax
    \@tempdimb=\f@size pt
    \hskip\VUespaceRatio\@tempdimb
      plus\VUespaceEtire\@tempdimb minus\VUespaceSerre\@tempdimb
  \else \hskip\spaceskip \fi}\makeatother
\begingroup
  \catcode`\ =\active
  \gdef\VUespacesInsecables{\catcode`\ =\active
    \def {\ifhmode\ifnum\lastnodetype>-1\relax\penalty10000\VUespaceMot\fi\fi}}%
\endgroup

% LA GRILLE DES LIGNES DE BASE EST IMPOSEE.
%
% TeX ne tient sa grille que tant que les lignes y entrent : si la
% hauteur naturelle d'une ligne depasse \baselineskip, il abandonne le
% pas regulier et separe les lignes par \lineskip, c'est-a-dire par
% leur encombrement reel. Le pas devient alors une consequence du
% texte, et non plus un reglage.
%
% C'est exactement ce qui se passait ici. Le fac-simile est compose
% presque au fer -- corps 10,93 pt sur un pas de 11,68 -- et le corps
% par page monte jusqu'a 12,43 pt pour remplir la justification. Au
% dela de ~11,7 pt, la ligne ne tient plus dans le pas, TeX lache la
% grille, et les pages s'allongent sans que rien ne le dise : vingt-neuf
% d'entre elles debordaient par le bas du papier, et changer
% l'interlignage de trois points n'y faisait presque rien -- la pente
% mesuree valait 0,5 la ou elle aurait du valoir 40.
%
% \lineskiplimit a -1000pt, la condition n'est jamais remplie et le pas
% reste celui qu'on a mesure, quel que soit le corps. C'est ce que fait
% le plomb : les lignes de base sont a distance fixe, imposee par le
% corps des interlignes, et une lettre un peu haute ne pousse pas sa
% voisine.
%
% LE REGLAGE SE POSE DANS LA PAGE, ET NON DANS LE PREAMBULE. Toute
% boite de LaTeX appelle \@parboxrestore, qui appelle \normalbaselines,
% qui remet \lineskiplimit a zero. Pose une fois pour toutes en tete du
% document, le reglage etait donc annule a l'ouverture de chaque
% minipage -- c'est-a-dire sur chacune des deux cent dix pages. On le
% repose dans l'environnement VUpage, apres l'ouverture de la boite.

% Aucune cesure decidee par TeX : toutes les coupures sont relevees.
\hyphenpenalty=10000
\exhyphenpenalty=10000
\pretolerance=-1
\tolerance=9999
\emergencystretch=0pt

% -------------------------------------------------------------------
% 3. PAGE FIGEE
% -------------------------------------------------------------------
% Chaque page du fac-simile est composee dans une minipage de hauteur
% fixe : une retouche dans un alinea ne peut pas faire migrer du texte
% vers la page suivante.
\pagestyle{empty}
\setlength{\parskip}{0pt}
\setlength{\topskip}{0pt}
\raggedbottom
\setlength{\parindent}{\VUalinea}

% LA MARQUE CACHEE, sur chaque page. Composee en MODE DE RENDU 3
% (« 3 Tr ») : le texte entre dans le flux du PDF, il est indexe,
% cherchable et extrait par pdftotext, mais il n'est PAS PEINT. Ce
% n'est pas du blanc sur blanc, qui se verrait en selectionnant : c'est
% du texte que le moteur de rendu saute. Corps 0,4 pt et boite de
% hauteur ET de largeur nulles, par surcroit de precaution.
\newcommand{\VUmarqueCachee}{%
  \vbox to 0pt{%
    \hbox to 0pt{%
      \pdfliteral{3 Tr}%
      \fontsize{0.4pt}{0.4pt}\selectfont\VUmarque
      \pdfliteral{0 Tr}\hss}%
    \vss}%
  \nointerlineskip}

% #1 (optionnel) = numero de FEUILLET du fac-simile. Il n'est pas
% compose : il sert aux controles a retrouver la page du scan.
% #2 = folio imprime ; vide si la page n'en porte pas.
% -------------------------------------------------------------------
% 3 bis. LE CORPS, PAGE PAR PAGE
% -------------------------------------------------------------------
% UN CORPS UNIQUE POUR TOUT LE VOLUME NE PEUT PAS REMPLIR TOUTES LES
% PAGES. Le corps global est celui qui, en moyenne, remplit le mieux ;
% il est donc, sur les pages ou l'original serre un peu plus, trop
% petit -- et les lignes s'y etirent, avec des blancs inter-mots
% beants. Le compositeur de 1926 ne tenait pas une chasse constante
% d'un cahier a l'autre, et la fonte d'aujourd'hui n'a pas exactement
% la sienne. On pose donc un corps PAR PAGE, mesure par
% outils/korpo.py : le plus grand qui ne fasse deborder aucune ligne.
\makeatletter
\newcommand{\VUkorpoPage}[2]{%
  \expandafter\gdef\csname VUkorpo@#1\endcsname{#2}}
\newcommand{\VUkorpoDe}[1]{%
  \ifcsname VUkorpo@#1\endcsname\csname VUkorpo@#1\endcsname\else
  \VUcorps\fi}
\newcommand{\VUtextePage}[1]{%
  \edef\VU@tmp{\VUkorpoDe{#1}}%
  \fontsize{\VU@tmp}{\VUinterlignoDe{#1}}\selectfont}
% L'INTERLIGNAGE AUSSI EST PROPRE A LA PAGE. Le pas mesure est une
% mediane ; les pages qui portent plus de lignes que la mediane ne
% tiennent pas dans le bloc et debordent par le bas. Le compositeur
% « cardait » ses pages -- il serrait ou ouvrait ses lames selon la
% matiere. outils/altezo.py mesure, page par page, le pas qui fait
% tenir la matiere dans la hauteur relevee.
\newcommand{\VUinterlignoPage}[2]{%
  \expandafter\gdef\csname VUintl@#1\endcsname{#2}}
\newcommand{\VUinterlignoDe}[1]{%
  \ifcsname VUintl@#1\endcsname\csname VUintl@#1\endcsname\else
  \VUinterligne\fi}
\makeatother

% LA PAGE SE COMPOSE DANS UNE BOITE QU'ON MESURE. La minipage a
% hauteur fixe se termine par un \vss, dont le retrait est INFINI :
% une page trop pleine ne produit donc aucun « Overfull \vbox », elle
% deborde en silence par le bas du papier. Quatre pages du livret ido
% et trois du francais sortaient ainsi du massicot sans que rien ne le
% signale, et c'est l'oeil qui les a trouvees.
%
% On compose donc d'abord dans une boite de hauteur naturelle, on
% annonce sa hauteur dans le journal — c'est outils/altezo.py qui la
% lit —, puis on la pose dans la page. Rien ne change a la composition ;
% ce qui change, c'est qu'on sait desormais ce qu'on compose.
\newsavebox{\VUpagoBoxo}
\newenvironment{VUpage}[2][]%
  {\par
   % LE \noindent VA AVANT LA BOITE POSEE, NON AVANT LA BOITE MESUREE.
   % Place ici, il ouvrait un alinea, la lrbox s'y inserait, et le vbox
   % de la page venait grossir cette meme ligne : cent onze pages sur
   % cent quatorze rendaient un « Overfull \hbox » de 11,3 pt, toujours
   % le meme, qui ne venait d'aucun texte.
   \begin{lrbox}{\VUpagoBoxo}%
   % LA HAUTEUR EST FIXE ICI, dans la boite mesuree elle-meme. Reglee
   % plus loin, au moment de poser la page, elle ramenait les 111
   % « Overfull \hbox » du renfoncement. Reglee ici, la boite sort deja
   % a la bonne hauteur et se pose sans enveloppe.
   % Contrepartie : \typeout{VUPAGO} rend alors toujours \textheight.
   % Pour mesurer les hauteurs naturelles, outils/altezo.py ote
   % temporairement le « [\textheight] » ci-dessous.
   \begin{minipage}[t][\textheight][t]{\textwidth}%
     % minipage remet \parindent a zero (\@parboxrestore) : on le
     % retablit, les deux fac-similes renfoncent tous leurs alineas.
     \setlength{\parindent}{\VUalinea}%
     \lineskiplimit=-1000pt \lineskip=0pt
     \VUespacesInsecables
     \VUnotofalse
     \VUtextePage{#1}\VUmarqueCachee\VUfolio{#2}\ignorespaces}%
  {\VUnotoPose
   \end{minipage}%
   \end{lrbox}%
   % HAUTEUR + PROFONDEUR. Une minipage alignee par le haut a son
   % point de reference sur la PREMIERE ligne de base : presque toute
   % sa matiere est donc en profondeur, et \ht seul rendait 0,0 pt
   % pour les cent quatorze pages.
   \typeout{VUPAGO \the\dimexpr\ht\VUpagoBoxo+\dp\VUpagoBoxo\relax}%
   % LA BOITE SE POSE TELLE QUELLE, sans etre reglee a \textheight.
   % Le \vbox to \textheight qui l'enveloppait ajoutait 11,3 pt a la
   % ligne qui la portait, sur chacune des deux cent dix pages -- un
   % « Overfull \hbox » par page, qui noyait les vrais. Il ne servait
   % qu'a empecher la matiere de migrer vers la page suivante, et le
   % \penalty-10000 qui suit l'empeche deja. La page commence de toute
   % facon a \VUmargeSup, que geometry fixe.
   % LE \noindent EST INDISPENSABLE. Sans lui la page recoit le
   % renfoncement d'alinea -- 11,29575 pt, soit exactement les 5 % de
   % la justification que vaut \VUalinea -- et chacune des deux cent
   % dix pages rendait un « Overfull \hbox » de cette largeur, toujours
   % la meme, qui ne venait d'aucun texte et noyait les vrais.
   % ET LA HAUTEUR RESTE FIXE. Posee a sa hauteur naturelle, une page
   % trop pleine se deversait sur une page de plus, et le volume en
   % gagnait une : la regle « une page du fac-simile = une page du
   % PDF » etait rompue. La boite est donc reglee a \textheight, ce qui
   % la garde a sa place ; les rares pages qui n'y tiennent pas sont
   % signalees par outils/altezo.py, et c'est leur releve qu'il faut
   % reprendre.
   \noindent\usebox{\VUpagoBoxo}%
   \par\penalty-10000\relax}

% Interlignage propre a une page — la « carde » du compositeur, qui
% repartit en lames uniformes le blanc d'une page trop courte. A
% appeler juste apres \begin{VUpage}, AVANT toute matiere. Sans appel,
% la page garde \VUinterligne.
\newcommand{\VUinterlignePage}[1]{\fontsize{\VUcorps}{#1}\selectfont}

% Blanc ORDINAIRE entre deux alineas. Il n'est PAS mis en \parskip : au
% sommet d'une minipage la colle de \parskip subsiste et decalerait
% toute la page d'autant. On l'ecrit donc, comme les sauts de ligne, la
% ou il est.
\newcommand{\VUblancAlinea}{\vskip\VUblancAlineaValeur\relax}
% Blanc plus large a une articulation. #1 = l'EXCES mesure sur le pas
% regulier de la page, DIMINUE du blanc ordinaire quand il le remplace.
\newcommand{\VUblanc}[1]{\vskip #1\relax}

% -------------------------------------------------------------------
% 4. CORPS DU TEXTE
% -------------------------------------------------------------------
\newcommand{\VUtexte}{\fontsize{\VUcorps}{\VUinterligne}\selectfont}
% Gras : famille propre, corps PROPORTIONNEL au corps courant — il suit
% donc automatiquement le texte, les notes, ou tout autre corps ou il
% apparait. Le gras des deux livrets est un demi-gras etroit qui n'est
% pas proportionne a son romain ; le gras de XCharter, lui, est large.
% D'ou une famille distincte et un rapport, non un corps absolu.
\makeatletter
\newcommand{\VUgras}[1]{{%
  \@tempdima=\f@size pt
  \@tempdima=\VUratioGras\@tempdima
  \fontfamily{\VUfonteGras}%
  \fontsize{\@tempdima}{\f@baselineskip}\selectfont\bfseries #1}}
\makeatother
\newcommand{\VUtexteNote}{%
  \fontsize{\VUcorpsNote}{\VUinterligneNote}\selectfont\lsstyle}

% -------------------------------------------------------------------
% 5. LIGNES D'APPARAT
% -------------------------------------------------------------------
% Sur ces pages aussi, UNE ligne du fac-simile = UN appel de macro : le
% releve reste ligne a ligne, et les controles peuvent le verifier.
% #1 = corps ; #2 = interlettrage (1/1000 de cadratin) ; #3 = texte.
% Les deux premiers sont MESURES : le corps par la hauteur de capitale,
% l'interlettrage par la largeur de la ligne (outils/apparat.py).
\newcommand{\VUcentre}[3]{{\centering\fontsize{#1}{#1}\selectfont
   \textls[#2]{#3}\par}}
\newcommand{\VUsaut}[1]{\vspace{#1}}

% Titre de section : capitales grasses espacees, centrees.
\newcommand{\VUtitre}[3]{%
  {\centering\fontfamily{\VUfonteGras}\fontsize{#1}{#1}\selectfont\bfseries
   \textls[#2]{#3}\par}}
% INTITULE EN PETITES CAPITALES — « DESKRIPTO GENERALA. », « LA BONA
% LERNANTI. », « I. DESCRIPTION GÉNÉRALE. ». Ni gras ni capitales
% pleines : c'est le romain du texte courant, dont seule la premiere
% lettre de chaque mot garde sa haute capitale. Les deux livrets en
% font le meme usage — c'est le rang de titre IMMEDIATEMENT SOUS celui
% du numero de tableau.
\newcommand{\VUintitulo}[2]{%
  {\centering\fontsize{#1}{#1}\selectfont\textls[#2]{\scshape}\par}}
\newcommand{\VUpk}[3]{%
  {\centering\fontsize{#1}{#1}\selectfont\textls[#2]{\scshape #3}\par}}

% Filet centre : #1 = largeur.
\newcommand{\VUfilet}[1]{{\centering\rule{#1}{0.4pt}\par}}
% Filet de note : court, cale a gauche.
\newcommand{\VUfiletnote}{{\noindent\rule{\VUfiletnoteLargeur}{0.4pt}\par}}

% Folio centre, entoure de deux demi-cadratins, comme « — 8 — ».
% Place ABSOLUMENT, en hauteur nulle : il ne pousse rien. Le corps de
% texte commence donc toujours exactement a \VUmargeSup, que la page
% porte un folio ou non.
\newcommand{\VUfolio}[1]{%
  \ifx\relax#1\relax\else
    \vbox to 0pt{\vskip\dimexpr\VUfolioY-\VUmargeSup\relax
      \hbox to \textwidth{\hss--\,#1\,--\hss}\vss}%
    \nointerlineskip
  \fi}

% -------------------------------------------------------------------
% 7 bis. NOTES DE BAS DE PAGE
% -------------------------------------------------------------------
% ELLES SE POSENT PAR LE BAS, ET NON PAR LE HAUT.
%
% Premiere version : chaque appel donnait l'ordonnee du FILET, et la
% consigne de releve proposait 143,2 mm « si la note est en bas d'une
% page pleine ». Les treize releves ont tous recopie cette valeur --
% c'est ma faute, pas la leur : une ordonnee de depart ne peut pas
% convenir a des notes de longueurs differentes. Une note de deux
% lignes flottait alors au milieu du blanc, et une note de douze
% lignes -- feuillet 19 du livret ido -- sortait par le bas du papier,
% ses trois dernieres lignes rognees au massicot.
%
% Le compositeur, lui, ne fait pas cela : il cale sa note SUR LE PIED
% du bloc, et c'est le haut du filet qui monte quand la note s'allonge.
% On mesure donc la note dans une boite, et on la pose de facon que son
% dernier rang tombe au bas du bloc de texte. L'ordonnee cesse d'etre
% une donnee du releve : elle se calcule, et elle est juste pour toutes
% les longueurs.
%
% #1 est CONSERVE mais ignore : le supprimer obligerait a retoucher les
% treize relevés sans rien y gagner, et l'argument garde trace de
% l'ordonnee que le fac-simile montre, si on veut un jour la verifier.
% ET LA NOTE NE SE POSE PAS EN HAUTEUR NULLE : ELLE SE DIFFERE.
%
% Deuxieme version : la note etait posee ABSOLUMENT, dans un \vbox to
% 0pt, cale sur le pied du bloc. Elle tombait juste -- et elle ne
% poussait rien. Sur une page dont le texte descend jusqu'au pied, les
% deux s'imprimaient donc l'un SUR l'autre : sept pages du livret ido
% et une du francais, mesurees a mutool en comparant le corps des
% lignes (la note vaut 0,79 fois le texte, l'exposant d'un renvoi 0,70
% -- ce sont les corps qui les separent, non la position).
%
% Le compositeur ne peut pas faire cela : sa note occupe le bas de la
% page, et le texte s'arrete au-dessus. On differe donc la note jusqu'a
% la FERMETURE de la page, ou elle entre dans le flux apres le dernier
% alinea. Le chevauchement devient impossible par construction, et non
% par un calcul qu'il faudrait refaire a chaque releve.
%
% LE RESSORT EST EN « fill » ET NON EN « fil » : la minipage de
% \VUpage porte deja un \vss au bas, et deux ressorts de meme ordre se
% partageraient le blanc -- la note se poserait au MILIEU de ce qui
% reste. Un ordre au-dessus prend tout, et la note descend au pied.
%
% #1 est CONSERVE mais ignore : le supprimer obligerait a retoucher les
% treize relevés sans rien y gagner, et l'argument garde trace de
% l'ordonnee que le fac-simile montre, si on veut un jour la verifier.
\newsavebox{\VUnotoBoxo}
\newsavebox{\VUnotoUno}
\newif\ifVUnoto
% Pied du bloc de texte, compte depuis le bord superieur du papier.
\newcommand{\VUnotoBas}{\dimexpr\VUfolioY+\VUtexteHauteur\relax}
% ET « \VUnotoBas » PARTAIT DU FOLIO, QUI EST AU-DESSUS DU BLOC.
% \VUfolioY donne la ligne de base du numero de feuillet, \VUmargeSup le
% sommet du bloc de texte : les deux different de 4,11 mm en ido et de
% 3,75 mm en francais. Caler la note sur \VUfolioY+\VUtexteHauteur la
% posait donc de cette hauteur-la TROP HAUT, c'est-a-dire dans le texte.
% Le pied du bloc est \VUmargeSup+\VUtexteHauteur, et c'est le bas de la
% minipage : la note s'y pose d'elle-meme, sans correction. \VUnotoBas
% n'est plus utilise que pour memoire.
\newcommand{\VUnotes}[2]{%
  % \vtop ET NON \vbox : le point de reference d'un \vbox est sa
  % DERNIERE ligne de base, de sorte qu'ajoute a la suite du texte il
  % remonte de toute sa hauteur et se pose PAR-DESSUS. Celui d'un \vtop
  % est sa premiere ligne : la note descend, comme on l'attend.
  \sbox{\VUnotoUno}{\vtop{\hsize=\textwidth
      \VUfiletnote\vskip\VUblancNote\VUtexteNote
      \setlength{\parindent}{\VUalineaNote}#2\par}}%
  % DEUX NOTES SUR UNE PAGE, cela arrive -- feuillet 41 du livret ido :
  % la seconde s'empile sous la premiere, elles ne se remplacent pas.
  \ifVUnoto
    \sbox{\VUnotoBoxo}{\vtop{\hbox{\usebox{\VUnotoBoxo}}\hbox{\usebox{\VUnotoUno}}}}%
  \else
    \sbox{\VUnotoBoxo}{\vtop{\hbox{\usebox{\VUnotoUno}}}}%
    \VUnototrue
  \fi
  \ignorespaces}
% Ce que \VUpage pose au bas de la page, s'il y a une note.
\newcommand{\VUnotoPose}{%
  \ifVUnoto
    \par\nobreak\vskip0pt plus1fill\relax
    \noindent\hbox to \textwidth{\usebox{\VUnotoBoxo}\hss}%
  \fi}

% -------------------------------------------------------------------
% 7 ter. CE QUI NE SE COMPOSE PAS
% -------------------------------------------------------------------
% Deux choses, dans ces deux volumes, ne sont pas de la composition :
% la gravure sur bois du tableau mural n^o 12, au feuillet 111 du
% livret ido, et l'encadrement Art nouveau de la quatrieme de
% couverture francaise. Les rendre en caracteres serait les inventer ;
% les omettre — ce qu'on avait fait d'abord, un \VUsaut et un
% commentaire — c'est publier une page qui ment sur ce qu'elle
% montrait. On les preleve donc sur le scan (outils/imaji.py) et on les
% pose a leur taille relevee.
%
% #1 = fichier ; #2 = largeur, en fraction de la justification.
\newcommand{\VUimajo}[2]{%
  {\centering\noindent
   \includegraphics[width=#2\textwidth]{ornamenti/#1}\par}}

% -------------------------------------------------------------------
% 7 quater. LIGNES D'AFFICHE ET ENCADREMENT
% -------------------------------------------------------------------
% \VUlineo : UNE LIGNE QUI NE SE JUSTIFIE PAS.
%
% Les pages d'annonce et la quatrieme de couverture sont faites de
% lignes courtes -- « 1^o En feuilles separees ; » -- que le releve
% notait comme du texte courant termine par \nl. Or \nl force une
% coupure DANS UN ALINEA JUSTIFIE : TeX etire alors la ligne d'une
% marge a l'autre, et « 1^o En feuilles separees ; » barrait toute la
% largeur en semant des blancs de deux cadratins entre ses mots. Le
% fac-simile, lui, pose ces lignes a leur chasse naturelle.
%
% Ces lignes ne sont donc pas du texte courant : ce sont des lignes
% d'affiche, et elles se composent comme telles -- au fer a gauche, a
% leur largeur propre, sans etirement.
% #1 = corps ; #2 = interlettrage ; #3 = retrait a gauche ; #4 = texte.
\newcommand{\VUlineo}[4]{%
  {\fontsize{#1}{#1}\selectfont
   \noindent\hspace*{#3}\textls[#2]{#4}\hfill\par}}
% Variante grasse, pour les egyptiennes d'affiche.
\newcommand{\VUlineoG}[4]{%
  {\fontfamily{\VUfonteGras}\fontsize{#1}{#1}\selectfont\bfseries
   \noindent\hspace*{#3}\textls[#2]{#4}\hfill\par}}

% \VUkadro : L'ENCADREMENT DE LA QUATRIEME DE COUVERTURE, en trace et
% non en image. Le fac-simile porte un cadre Art nouveau : un double
% filet dont les angles s'infléchissent en volutes. On ne le photo\-
% graphie pas — ce serait renoncer a composer la page — on le dessine.
% Le trace ci-dessous n'est pas une copie du bois original : c'est un
% cadre de meme facture, double filet a angles adoucis et fleurons
% d'angle, pose aux dimensions relevees. La difference est dite ici
% plutot que masquee.
% #1 = hauteur du cadre.
\newcommand{\VUkadro}[2]{%
  \begingroup
  \setlength{\unitlength}{1pt}%
  \noindent\begin{tikzpicture}[baseline=(courant.north)]
    % inner sep=0pt : sans lui, ce noeud vide reste cadre par la marge
    % interne par defaut (0,333 em) et elargit la boite de 3,88 pt au
    % dela de \textwidth -- un « Overfull \hbox » constant, identique
    % sur toutes les pages, quel que soit leur contenu.
    \node[inner sep=0pt] (courant) at (0,0) {};
    % Le filet exterieur est trace 0,55 pt EN DEDANS de chaque bord : un
    % trait de 1,1 pt centre sur x=0 deborde de 0,55 pt a gauche ET a
    % droite de \textwidth, et ce dernier debord rendait un « Overfull
    % \hbox » de 1,1 pt sur la page entiere.
    \draw[line width=1.1pt, rounded corners=7pt]
      (0.55pt,0) rectangle (\textwidth-0.55pt, -#1);
    \draw[line width=0.4pt, rounded corners=5pt]
      (4pt,-4pt) rectangle (\textwidth-4pt, -#1+4pt);
    % Le decalage vertical se calcule A PART, par \pgfmathsetmacro : passe
    % directement dans la clef shift={...}, pgfmath ne l'evaluait pas --
    % les quatre fleurons se posaient tous pres du haut, deux dedans, deux
    % dehors, et le bas du cadre restait nu.
    \foreach \x/\y/\rx/\ry in {0/0/1/-1, 1/0/-1/-1, 0/1/1/1, 1/1/-1/1} {
      \pgfmathsetmacro{\VUkadroY}{-\y*0 - (1-\y)*#1}
      \begin{scope}[shift={(\x*\textwidth, \VUkadroY pt)}]
        \draw[line width=0.9pt]
          (\rx*13pt, \ry*3pt) .. controls (\rx*7pt, \ry*10pt)
          and (\rx*3pt, \ry*7pt) .. (\rx*3pt, \ry*13pt);
      \end{scope}}
    \node[inner sep=0pt, anchor=north west] at (11pt, -13pt)
      {\begin{minipage}{\dimexpr\textwidth-22pt\relax}#2\end{minipage}};
  \end{tikzpicture}\par
  \endgroup}

% Ornement typographique (le fleuron des pages de tableau). Le trace du
% fac-simile n'est pas de la fonte : il se pose en image quand on l'a
% releve, et se remplace par un tilde espace tant qu'on ne l'a pas.
\newcommand{\VUornamento}[1]{{\centering\fontsize{#1}{#1}\selectfont
  \textls[300]{\textasciitilde\textasciitilde\textasciitilde}\par}}

% -------------------------------------------------------------------
% 8. LA TABLE DES CORPS PAR PAGE
% -------------------------------------------------------------------
% EN DERNIER, ET C'EST LA TOUT. Ce chargement etait d'abord place au
% debut du fichier, avant meme \geometry. \VUkorpoPage n'y etait pas
% encore definie : LaTeX prenait donc chaque appel pour du texte, et
% composait « 911.60pt 1011.15pt 1111.30pt... » -- le contenu du
% fichier de reglage imprime au fil des pages. Le symptome etait dix
% « Overfull \hbox » inexplicables ; la cause etait un ordre.
% Le fichier peut ne pas exister : le volume se compose alors au corps
% global, ce qui est le comportement voulu avant la premiere mesure.
\InputIfFileExists{korpi-\VUlangue.tex}{}{}
